% === Slide: <<NAME>> ===
% LAYOUT: code-window — a single centred monospace box (prompt, code, I/O
% example) over a full-bleed image (typically a "window"/terminal frame image).
% Box: centre (0.50pw, 0.55ph), text width 0.61pw, minimum height 0.66ph,
% opacity 0.82 (most opaque card: text legibility first).
% Token colouring convention inside the box:
%   \textcolor{NoirAccent}{\bfseries Label:}  section labels (Question:, Answer:)
%   \textcolor{NoirCyan}{<token>}             variables / placeholders / outputs
%   \textcolor{NoirMuted}{...}                 de-emphasised text
% Escape _ as \_ and use \\[2pt] between lines, \\[5pt]/\\[8pt] between sections.
% Use with \NoirBackgroundVoid instead if no window image is available.
\NoirBackgroundImage{<<IMAGE_FILE>>}
\begin{frame}{<<FRAME_TITLE>>}
\begin{tikzpicture}[remember picture, overlay,
    codebox/.style={
        fill=NoirBg, fill opacity=0.82, text opacity=1,
        rounded corners=6pt,
        text width=0.61\paperwidth,
        align=left, inner sep=16pt,
        anchor=center,
        minimum height=0.66\paperheight
    }
]
  \node[codebox] at ([xshift=0.50\paperwidth, yshift=-0.55\paperheight] current page.north west) {%
    {\scriptsize\color{NoirText}\ttfamily
      \textcolor{NoirAccent}{\bfseries <<SECTION_1_LABEL>>}\\[5pt]
      <<SECTION_1_LINES>>\\[8pt]
      \textcolor{NoirAccent}{\bfseries <<SECTION_2_LABEL>>}\\[5pt]
      <<SECTION_2_LINES>>%
    }%
  };
\end{tikzpicture}
\end{frame}
\NoirBackgroundReset
